% !TeX program = xelatex
% Compilation : xelatex rapport_esiea.tex (2 passes pour le sommaire/numérotation)
% Nécessite XeLaTeX ou LuaLaTeX (fontspec) — pdflatex ne fonctionnera pas.
%
% Polices requises (à placer dans un dossier fonts/ à côté de ce fichier,
% téléchargeables gratuitement sur fonts.google.com — licence Apache) :
%   fonts/Poppins-Bold.ttf, Poppins-SemiBold.ttf, Poppins-Regular.ttf
%   fonts/Roboto-Regular.ttf, Roboto-Bold.ttf, Roboto-Italic.ttf
% Si les polices ne sont pas installées système, décommentez la ligne
% \setmainfont{Roboto}[Path=fonts/,...] plus bas avec le chemin explicite.

\documentclass[11pt,a4paper]{report}

% ---------- Encodage / langue ----------
\usepackage{fontspec}
\usepackage[french]{babel}
\usepackage{microtype}

% ---------- Couleurs charte ESIEA ----------
\usepackage[table]{xcolor}
\definecolor{esieablue}{HTML}{36A9E1}      % Bleu clair — couleur principale
\definecolor{esienavy}{HTML}{151D34}       % Bleu marine — couleur secondaire foncée
\definecolor{esieaorange}{HTML}{F59E00}    % Orange — accent (barres, encarts)
\definecolor{esieaoffwhite}{HTML}{FDF8F1}  % Blanc cassé — fonds clairs
\definecolor{esieagray}{HTML}{4A4A4A}      % Gris texte courant (dérivé, lisibilité)

% ---------- Polices (Poppins = titres, Roboto = corps) ----------
% Si Poppins/Roboto sont installées au niveau système, ces deux lignes suffisent :
\setmainfont{Roboto}
\newfontfamily\titlefont{Poppins}[BoldFont={Poppins SemiBold}]
% Sinon, remplacez par (chemins relatifs au fichier .tex) :
% \setmainfont{Roboto-Regular.ttf}[Path=fonts/,BoldFont=Roboto-Bold.ttf,ItalicFont=Roboto-Italic.ttf]
% \newfontfamily\titlefont{Poppins-SemiBold.ttf}[Path=fonts/,BoldFont=Poppins-Bold.ttf]

% ---------- Mise en page ----------
\usepackage[margin=2.5cm,top=3cm]{geometry}
\usepackage{titlesec}
\usepackage{titletoc}
\usepackage{fancyhdr}
\usepackage{graphicx}
\usepackage{tikz}
\usepackage{tcolorbox}
\tcbuselibrary{skins,breakable}
\usepackage{booktabs}
\usepackage{longtable}
\usepackage{hyperref}
\usepackage{enumitem}
\usepackage{caption}

% ---------- Titres de section façon charte (barre orange sous le titre) ----------
\titleformat{\chapter}[display]
  {\titlefont\Huge\bfseries\color{esienavy}}
  {\filright\normalfont\normalsize\color{esieablue}\MakeUppercase{Rapport Big Data — IA 2025-2026}}
  {6pt}
  {\titlerule[2pt]\color{esieaorange}\vspace{4pt}\thechapter.\ }
[\vspace{2pt}{\color{esieaorange}\titlerule[3pt]}]

\titleformat{\section}
  {\titlefont\Large\bfseries\color{esieablue}}{\thesection}{10pt}{}
\titlespacing{\section}{0pt}{18pt}{8pt}

\titleformat{\subsection}
  {\titlefont\large\bfseries\color{esienavy}}{\thesubsection}{8pt}{}

% ---------- En-têtes / pieds de page ----------
\pagestyle{fancy}
\fancyhf{}
\renewcommand{\headrulewidth}{0.6pt}
\renewcommand{\footrulewidth}{0pt}
\fancyhead[L]{\small\color{esienavy}\textbf{esiea} \,|\, Projet Big Data IA}
\fancyhead[R]{\small\color{esieablue}\leftmark}
\fancyfoot[C]{\small\color{esieagray}\thepage}

% ---------- Encadré "constat / interprétation" réutilisable ----------
\newtcolorbox{constat}[1][]{
  colback=esieaoffwhite, colframe=esieablue, boxrule=0.8pt,
  arc=2pt, left=8pt, right=8pt, top=6pt, bottom=6pt,
  fonttitle=\titlefont\bfseries\color{white}, coltitle=white,
  colbacktitle=esieablue, title={#1}, breakable
}
\newtcolorbox{recommandation}[1][]{
  colback=white, colframe=esieaorange, boxrule=1pt,
  arc=2pt, left=8pt, right=8pt, top=6pt, bottom=6pt,
  fonttitle=\titlefont\bfseries\color{esienavy}, title={#1}, breakable
}

\hypersetup{
  colorlinks=true, linkcolor=esieablue, citecolor=esieablue,
  urlcolor=esieablue, pdftitle={Marché de l'emploi en IA 2025-2026},
  pdfauthor={Équipe projet — ESIEA}
}

\title{\titlefont\Huge\bfseries\color{esienavy}Marché de l'emploi en IA\\[4pt]\Large 2025--2026}
\author{}
\date{}

% =========================================================
\begin{document}

% ---------- Page de garde façon charte (fond marine, titre blanc, barre orange) ----------
\begin{titlepage}
\newgeometry{margin=0cm}
\begin{tikzpicture}[remember picture, overlay]
  \fill[esienavy] (current page.south west) rectangle (current page.north east);
\end{tikzpicture}
\vspace*{3cm}
\begin{center}
  {\titlefont\color{white}\Large esiea}\\[2.5cm]
  {\titlefont\color{white}\Huge\bfseries Projet Big Data\\[6pt] Marché de l'emploi en IA}\\[10pt]
  {\titlefont\color{white}\huge 2025 -- 2026}\\[18pt]
  {\color{esieaorange}\rule{6cm}{3pt}}\\[18pt]
  {\color{white}\large HDFS $\cdot$ Hive $\cdot$ Spark $\cdot$ Jupyter/PySpark}\\[40pt]
  {\color{white}\normalsize Sujet proposé par M.~Touati Zouhaier}\\[4pt]
  {\color{white}\normalsize Équipe : \emph{[Prénom NOM $\times$5]}}\\[4pt]
  {\color{white}\normalsize ESIEA — Spécialisation DevSecOps \& Cloud Computing}\\[4pt]
  {\color{white}\normalsize Septembre 2026}
\end{center}
\end{titlepage}
\restoregeometry

\tableofcontents
\clearpage

% ---------- Nécessite \usepackage{tikz} pour la page de garde ----------

\chapter{Introduction}
\section{Contexte}
% Présenter le contexte du marché de l'emploi en IA.

\section{Problématique}
% Formuler la problématique du projet.

\section{Objectifs}
% Objectifs techniques et analytiques.

\section{Intérêt de la chaîne HDFS $\to$ Hive $\to$ Spark $\to$ Jupyter/PySpark}
% Expliquer brièvement l'intérêt de l'architecture retenue.

\chapter{Présentation du dataset}
% Source, nombre de lignes/colonnes, période couverte, variables principales,
% premières observations sur la qualité des données.

\chapter{Architecture Big Data retenue}
\section{Schéma d'architecture}
% \includegraphics[width=\linewidth]{figures/architecture.png}

\section{Rôle de chaque composant}
% HDFS, Hive, Spark, Jupyter/PySpark (remplace Zeppelin — cf. note ci-dessous).

\begin{constat}[Choix technique validé]
Choix validé avec M.~Touati~: Jupyter/PySpark remplace Zeppelin comme
interface d'analyse et de dashboard. HDFS et Hive restent le socle de
stockage et de requêtage, conformément au sujet.
\end{constat}

\section{Organisation raw / processed / output}

\chapter{Ingestion des données dans HDFS}
% Arborescence HDFS, commandes mkdir/put/ls/du, vérification du nombre de lignes.

\chapter{Création et contrôle de la table Hive RAW}
% Schéma retenu, CREATE EXTERNAL TABLE, DESCRIBE, SELECT LIMIT, COUNT(*).

\chapter{Qualité et préparation des données avec Spark}
% Chargement, printSchema(), contrôle NULL/doublons/types, règles de nettoyage justifiées.

\chapter{Construction des tables offres, metiers et competences}
% Grain de chaque table, clés job\_id/metier\_id, split()/explode(), schéma relationnel.

\chapter{Analyses Hive SQL et Spark SQL}
% Pour chaque analyse : question -> tables -> requête -> résultat -> interprétation.

\chapter{Analyse des métiers, salaires et compétences}

\chapter{Dashboard Jupyter/PySpark}
% Minimum 6 visualisations, titre + question + interprétation chiffrée sous chaque graphique.

\chapter{Synthèse des résultats}
\begin{constat}[Constat 1]
% Relier chaque constat à un indicateur ou une requête précise.
\end{constat}

\chapter{Recommandations}
\begin{recommandation}[Recommandation 1]
% Résultat observé -> interprétation -> recommandation.
\end{recommandation}

\chapter{Limites de l'étude}
% Représentativité, volume du dataset, limites techniques/méthodologiques.
% Rappeler que le dataset ne mesure pas l'impact causal de l'IA sur l'emploi global.

\chapter{Conclusion}

\appendix
\chapter{Annexes}
\section{Commandes HDFS}
\section{Scripts PySpark}
\section{CREATE TABLE Hive}
\section{Requêtes Hive SQL / Spark SQL}
\section{Journal d'exécution}
% Voir JOURNAL.md du dépôt — synthèse à intégrer ici avant le rendu final.

\end{document}
